% ---------------------------------------------------------------
% Section 4 — Implementation in the BX3 Framework
% Paper: Bailout Protocol: Mandatory Human Accountability in Multi-Agent Systems
% ---------------------------------------------------------------

\section{Implementation in the BX3 Framework}
\label{sec:bailout-implementation}

The Bailout Protocol is not a safety module bolted onto an existing system.
It is the \emph{runtime expression} of the BX3 Framework's upstream accountability
guarantee --- the mechanism that converts a declarative layer commitment into a
binding, non-configurable, auditable operational decision.  The protocol integrates
with each layer of the BX3 architecture: the Purpose Layer provides the human
accountability anchor; the Bounds Engine serves as the exclusive triggering
mechanism; the Fact Layer enforces the resulting audit trail; and the Forensic
Ledger provides cryptographic evidence that the guarantee was honoured.  This
section details each integration point, the state machine that governs every
trigger's lifecycle, and the formal relationship between the protocol and BX3's
upstream intent.

% ----------------------------------------------------------------
\subsection{Purpose Layer as Human Accountability Anchor}
% ----------------------------------------------------------------

The Purpose Layer (\texttt{PurposeLayer.ts}) defines the \emph{intent, judgment,
and accountability boundary} for every node in the BX3 execution tree.  Each
node is initialised with three canonical, immutable fields:

\begin{itemize}
  \item \texttt{nodeId} --- a unique identifier for this agent or process.
  \item \texttt{humanRoot} --- the designated human decision-maker at the
        terminal boundary of the accountability chain for this node's subtree.
        This field is the protocol's destination: propagation terminates here.
  \item \texttt{accountabilityAnchor} --- a stable, versioned reference string
        that identifies the governing policy regime under which this node operates.
        This enables post-hoc reconstruction of which policy was active at
        firing time, independent of subsequent policy revisions.
\end{itemize}

These fields are sealed at node initialisation and propagated downward through
the execution tree.  When a child node is spawned, it receives a \texttt{parentId}
and a \texttt{parentPurposeHash} that cryptographically binds it to its parent's
purpose.  The chain is tamper-evident: any divergence between a child's purpose
and its parent's is detectable by recomputing the parent's hash and comparing it
to the stored value.

The Bailout Protocol uses the Purpose Layer at two distinct stages.  First, at
trigger time, it resolves the \texttt{humanRoot} by traversing the
\texttt{propagationChain} upward until it finds the node whose \texttt{nodeId}
matches \texttt{humanRoot}.  This lookup is deterministic and bounded by tree
depth.  No autonomous resolution is permitted beyond this node by construction
of the state machine (Section~\ref{sec:state-machine}).  Second, at fire time,
the current \texttt{accountabilityAnchor} is written into the trigger's
\texttt{event} payload, permanently attaching the governing policy context to
the bail-out record.  A policy revision after the fact cannot rewrite the anchor
that was recorded.

The Purpose Layer thus provides the \emph{semantic grounding} for every
bail-out: it answers not only ``who is responsible?'' but also ``under which
governing commitment was this escalation mandatory?''  This distinction is
critical for regulatory defensibility, where the complete policy context
surrounding a trigger must be reconstructible, not merely the operational chain
of nodes.

% ----------------------------------------------------------------
\subsection{Bounds Engine: The Triggering Mechanism}
% ----------------------------------------------------------------

The Bounds Engine (\texttt{BoundsEngine.ts}) implements BX3's bounded-reasoning
layer, responsible for generating decision proposals (P5) and projecting their
consequences with confidence scores (P6).  The Bounds Engine is the
\textbf{exclusive source} of bail-out trigger events in the BX3 runtime ---
no other component may directly invoke \texttt{BailoutProtocol.fire()}.  This
exclusivity is architecturally enforced: the Bailout Protocol's internal API
accepts trigger requests only from the Bounds Engine, and the dependency is
checked at module initialisation.

Triggering occurs when the Bounds Engine evaluates a proposal against the
active \texttt{SafetyEnvelope} and detects one of three named
\texttt{TriggerCondition} types:

\begin{verbatim}
type TriggerCondition =
  | 'capability_boundary'
      // proposed action exceeds this node's defined capability set
  | 'safety_envelope_predicted'
      // P6 projection forecasts a SafetyEnvelope violation at t + delta
  | 'accountability_boundary'
      // proposed action requires human authorisation by active policy
\end{verbatim}

Each condition is evaluated against domain-specific constraints stored in the
\texttt{SafetyEnvelope} struct (e.g., water-rights limits, soil-moisture
thresholds, equipment load ratings).  The Bounds Engine does not block actions
directly: upon detecting a condition, it generates a \texttt{BailoutTrigger}
record and returns control to the Bailout Protocol orchestrator.  This
separation is architectural.  The Bounds Engine \emph{reasons about possibility};
the Bailout Protocol \emph{governs what happens when possibility exceeds
permission}.

When a trigger condition is met, the protocol calls \texttt{fire()}:

\begin{verbatim}
fire({
  triggerCondition: TriggerCondition,
  confidence: number,           // P6 confidence score
  event: Record<string,unknown>,
  nodeState: Record<string,unknown>,
  proposedResolution: string | null
}): BailoutTrigger
\end{verbatim}

The returned \texttt{BailoutTrigger} carries a unique \texttt{triggerId}, the
\texttt{originatingNodeId}, the full \texttt{propagationChain}, and the
P6 \texttt{confidence} score.  The \texttt{propagationChain} is initialised
with the originating node at this point.  The trigger is immediately written
to the Forensic Ledger as a \texttt{TRIGGER\_FIRED} entry before any further
processing occurs.  This write-before-propagate ordering is load-bearing: it
ensures that the trigger event is attested in the ledger even if a later
intermediate node is compromised and fails to propagate.

The \texttt{confidence} field from the Bounds Engine's P6 projection is
recorded but does not filter or delay the trigger.  A trigger fires regardless
of confidence level; the confidence score is evidence metadata, not a
gate.  The trigger condition type alone determines that the human must review.

% ----------------------------------------------------------------
\subsection{Fact Layer: Enforcement Partner and Audit Trail}
% ----------------------------------------------------------------

The Fact Layer (\texttt{FactLayer.ts}) is BX3's runtime enforcement layer ---
the component that evaluates proposals against observed reality and issues the
verdict that permits or blocks execution.  It serves as the Bailout Protocol's
enforcement partner in a two-key architecture: the Fact Layer clears the
technical path; the Bailout Protocol clears the accountability path.  Both
must be satisfied before any autonomous action proceeds to physical execution.

The Fact Layer's \texttt{evaluate()} method --- the Sandbox Gate ---
is the architectural chokepoint for all P5-to-P8 transitions.  It accepts a
\texttt{P5Decision} and its corresponding P6 projection and returns one of
three verdicts:

\begin{verbatim}
type SandboxVerdict = 'EXECUTE' | 'BLOCK' | 'REVIEW'
\end{verbatim}

The \texttt{EXECUTE} verdict is necessary but not sufficient for actionuation.
Even after a proposal passes the Sandbox Gate, the Bailout Protocol may suppress
execution if the associated trigger is in \texttt{human\_review} status.  The
state machine enforces this: execution is unblocked only after
\texttt{status} reaches \texttt{resolved} (P8) and the \texttt{RESOLUTION\_RECORDED}
ledger entry has been written.  If the trigger is in any other status, the Fact
Layer's \texttt{EXECUTE} verdict is held --- the action proceeds only when both
layers have independently cleared it.

Conversely, the Fact Layer enforces its own audit obligations.  Each
verdict issued by \texttt{evaluate()} writes a corresponding
\texttt{FACT\_LAYER\_VERDICT} entry to the Forensic Ledger:

\begin{enumerate}
  \item \texttt{FACT\_LAYER\_VERDICT(EXECUTE)} --- logged when the Fact Layer
        clears a proposal.  The corresponding action may proceed to the
        Sandbox Engine for execution unless the Bailout Protocol hold is active.
  \item \texttt{FACT\_LAYER\_VERDICT(BLOCK)} --- logged when the Fact Layer
        blocks a proposal due to detected reality violation.  The trigger
        may fire independently from the Bailout Protocol if a condition is met.
  \item \texttt{FACT\_LAYER\_VERDICT(REVIEW)} --- logged when the Fact Layer
        requires human review before any verdict is issued.  The
        \texttt{REVIEW} status blocks execution until the human provides
        a verdict; this state is coordinated with the Bailout Protocol's
        \texttt{human\_review} state through a shared \texttt{triggerId}.
  \item \texttt{ACTUATION} --- logged when the Sandbox Engine executes an
        action.  This is the terminal event that confirms the full
        P1--P8 chain has completed.  The \texttt{triggerId} links this entry
        to the originating trigger, enabling full path reconstruction from
        fire to actuation.
\end{enumerate}

The Fact Layer thus provides two functions simultaneously: it is the
technical gatekeeper that evaluates proposals against observed state, and it
is the co-enforcer of the accountability chain that records every gate passage
in the Forensic Ledger.  No decision passes without a ledger entry; no ledger
entry exists without a corresponding decision.  This mutual entailment is what
makes the audit trail self-consistent and resistant to retrospective fabrication.

% ----------------------------------------------------------------
\subsection{Forensic Ledger: Cryptographic Event Recording}
% ----------------------------------------------------------------

The Forensic Ledger (\texttt{ForensicLedger.ts}) is a write-once append-only
log that records every state transition, trigger event, and human resolution
in the Bailout Protocol.  It is the evidentiary substrate of BX3's accountability
guarantee --- the artifact that renders the upstream commitment machine-verifiable
and auditor-accessible without requiring access to the runtime system itself.

The ledger is structured as a chain of \texttt{LedgerEntry} records:

\begin{verbatim}
interface LedgerEntry {
  index: number;            // position in chain, 0-indexed
  timestamp: string;        // ISO 8601, UTC
  eventType: string;        // e.g., 'TRIGGER_FIRED', 'ESCALATED'
  triggerId: string;        // links to BailoutTrigger.triggerId
  nodeId: string;           // the node that authored this entry
  payload: Record<string,unknown>;  // event-specific data
  previousSeal: string;     // SHA-256 seal of the preceding entry
  seal: string;             // SHA-256(payload || previousSeal || timestamp)
}
\end{verbatim}

The \texttt{seal} is computed as:

\begin{equation}
\label{eq:ledger-seal}
\text{seal} = \text{SHA256}\bigl(
  \text{JSON.stringify}(\text{payload}) \;\|\; \text{previousSeal} \;\|\; \text{timestamp}
\bigr)
\end{equation}

where \texttt{previousSeal} is the \texttt{seal} of the immediately preceding
entry, or the string literal \texttt{'GENESIS'} for the first entry.  This
construction is isomorphic to a blockchain block header: any modification of a
historical entry invalidates all subsequent seals, because each subsequent
entry's recomputed seal will not match the stored value.  An adversary who
alters entry $k$ must re-compute and rewrite entries $k+1, k+2, \ldots, n$
to maintain consistency --- a cost linear in chain length.

The ledger exposes two operations:

\begin{itemize}
  \item \texttt{append(params)} --- atomically appends a new entry, capturing
        the previous seal at write time.  The ledger is thus equivalent to a
        write-once register: an entry, once written, cannot be altered without
        detection.
  \item \texttt{verify()} --- performs an $O(n)$ full-chain traversal that
        recomputes each entry's seal from its payload, previous seal, and
        timestamp and compares the result to the stored seal.  Returns
        \texttt{valid: true} and the chain length if all seals match; returns
        \texttt{valid: false} and the index of the first corrupted entry if any
        inconsistency is detected.
\end{itemize}

For the Bailout Protocol specifically, the ledger records four event types:

\begin{enumerate}
  \item \texttt{TRIGGER\_FIRED} --- logged at P1 when \texttt{fire()} is called,
        capturing the full \texttt{triggerCondition}, \texttt{confidence}, and
        \texttt{nodeState} snapshot.  Written before propagation begins.
  \item \texttt{ESCALATED} --- logged at each P3--P5 hop, recording the
        intermediate \texttt{nodeId} and propagation direction.  A gap in this
        sequence (an expected node ID absent from its temporal position) is
        evidence of a failed or suppressed propagation.
  \item \texttt{HUMAN\_REVIEW} --- logged at P7 when the trigger arrives at the
        \texttt{humanRoot}, recording the arrival timestamp and the pending
        decision.  The human's identity is verifiable through the
        \texttt{humanRoot} field paired with external authentication logs.
  \item \texttt{RESOLUTION\_RECORDED} --- logged at P8 when the human resolves
        the trigger, recording the resolution text, the resolver's identity,
        and the end-to-end latency from \texttt{firedAt} to \texttt{resolvedAt}.
        This is the terminal entry for the trigger's lifecycle.
  \item \texttt{TRIGGER\_EXPIRED} --- logged at P9 if the TTL expires before
        resolution, recording the deadline and the state at time of expiry.
\end{enumerate}

The combination of the \texttt{propagationChain} array inside the
\texttt{BailoutTrigger} record and the corresponding \texttt{LedgerEntry}
sequence in the Forensic Ledger provides dual-mode auditability: the chain
gives a human-readable provenance graph; the ledger gives a machine-verifiable
cryptographic proof of integrity.  Both are generated from the same runtime
events, eliminating the possibility of divergence between operational records
and forensic records.

% ----------------------------------------------------------------
\subsection{State Machine as Fact Layer Extension}
\label{sec:state-machine}
% ----------------------------------------------------------------

The Bailout Protocol is implemented as a deterministic finite state machine
(DFSM) whose states are a superset of the BX3 layer-plane taxonomy.
The state machine governs the lifecycle of every trigger and is the mechanism
through which the Fact Layer's two-key enforcement architecture is realised
in time.

The protocol defines five states, each mapped to a \texttt{status} field in
the \texttt{BailoutTrigger} record:

\bigskip
\noindent
\begin{tabular}{llp{8cm}}
  \hline
  \textbf{State} & \textbf{Plane} & \textbf{Semantics} \\
  \hline
  \texttt{pending}       & P1 & Trigger authored; propagation not yet begun \\
  \texttt{escalating}   & P3--P5 & Propagating upward through intermediate nodes \\
  \texttt{human\_review} & P7 & Arrived at \texttt{humanRoot}; human is reviewing \\
  \texttt{resolved}     & P8 & Human has issued a binding resolution; execution unblocked \\
  \texttt{expired}      & P9 & TTL exceeded without human resolution; failure state \\
  \hline
\end{tabular}
\bigskip

State transitions are driven by three primitive operations:

\begin{itemize}
  \item \texttt{fire()} --- P1 transition.  Creates the \texttt{BailoutTrigger}
        record, initialises \texttt{propagationChain} with the originating node,
        sets \texttt{status} to \texttt{pending}, and appends \texttt{TRIGGER\_FIRED}
        to the Forensic Ledger.

  \item \texttt{receiveAndPropagate()} --- P3--P5 transitions.  Each intermediate
        node receives the trigger, appends its \texttt{nodeId}, action
        (\texttt{ESCALATED}), and timestamp to \texttt{propagationChain}, advances
        \texttt{status} to \texttt{escalating}, and forwards the trigger toward the
        next node in the chain.  A corresponding \texttt{ESCALATED} ledger entry
        is written at each hop.

  \item \texttt{resolve(\{resolutionText, resolverId\})} --- P7/P8 transition.
        The \texttt{humanRoot} node receives the trigger, changes \texttt{status}
        to \texttt{human\_review} (P7), records the human's decision, and advances
        \texttt{status} to \texttt{resolved} (P8).  A \texttt{RESOLUTION\_RECORDED}
        ledger entry is appended.  Execution is unblocked only after this
        transition completes; the Fact Layer holds the action until it observes
        the \texttt{resolved} status.
\end{itemize}

The \texttt{expired} state is a defined failure state.  If the TTL expires before
a \texttt{resolve()} call, the trigger enters \texttt{expired} status and the
Bailout Protocol emits a \texttt{TRIGGER\_EXPIRED} ledger entry.  The specific
failure policy --- halt, default action, degraded mode --- is a domain-specific
parameter passed to the protocol at initialisation and is outside the protocol's
own scope.  Critically, the \texttt{expired} state does not authorise autonomous
continuation: the underlying condition remains ungoverned by the human, and the
system enters a distinct failure mode that requires operational intervention.
This preserves the safety guarantee even when the liveness guarantee is violated.

The state machine's determinism is load-bearing for the audit system.  Given the
same trigger event and the same \texttt{propagationChain}, the protocol always
produces the same sequence of state transitions and the same ledger entries.
This repeatability is what allows the Forensic Ledger to serve as authoritative
evidence in audits and post-mortems: the causal chain from trigger to resolution
is uniquely determined by the initial conditions and the state machine's
transition function, and that chain is entirely reconstructible from the ledger.

The state machine is also the mechanism by which the Fact Layer's two-key
architecture acquires temporal dimension.  The Fact Layer evaluates at a single
instant; the state machine extends that evaluation across the full lifecycle of
the trigger.  If the Fact Layer issues \texttt{EXECUTE} but the trigger is still
in \texttt{escalating} status, execution is held.  If the trigger reaches
\texttt{human\_review}, the Fact Layer's prior verdict remains valid but the
accountability path is not yet cleared.  The state machine coordinates these
two independent enforcement mechanisms across time, ensuring that neither
technical safety nor human accountability can be bypassed by either party acting
alone.

% ----------------------------------------------------------------
\subsection{Relationship: Bailout Protocol as Runtime Expression of
            the Upstream Accountability Guarantee}
% ----------------------------------------------------------------

The BX3 Framework's accountability guarantee is formally expressed in the
Intent Layer's constraint: \emph{for every intent $i \in I$, there exists a
role $r \in R$ and an action $a \in A$ such that $i \rightarrow (r, a)$}.
This guarantee is declarative --- it describes what must be true about the
system's design.  The Bailout Protocol is the mechanism by which this
declaration becomes operative at runtime.

The causal chain from declaration to enforcement is:

\begin{enumerate}
  \item The \texttt{humanRoot} field in the Purpose Layer \emph{declares} that a
        specific human is the final accountable authority for actions originating
        under this node's governance subtree.
  \item The \texttt{accountability\_boundary} condition in the Bounds Engine
        \emph{detects} when a proposed action falls outside the scope of
        autonomous authorisation --- when policy requires human decision regardless
        of technical feasibility.
  \item The Bailout Protocol's \texttt{fire()} method \emph{operationalises} the
        detection by creating a trigger that is obligated to propagate to the
        declared human.  Propagation is not discretionary; the state machine
        enforces it.
  \item The Fact Layer \emph{clears the technical path} with \texttt{EXECUTE},
        but the state machine \emph{blocks the accountability path} until
        \texttt{status} reaches \texttt{resolved}.  Both must be satisfied.
  \item The Forensic Ledger \emph{records} every step of this chain, providing
        cryptographic evidence that the guarantee was honoured in practice ---
        not merely declared in documentation.
\end{enumerate}

If any step in this chain is absent --- if the Bounds Engine fails to detect
the condition, if the protocol fails to propagate, if the state machine permits
autonomous override, or if the ledger is not written --- the upstream
accountability guarantee is voided.  The Bailout Protocol is therefore not an
additional feature layered on top of BX3; it is the \emph{necessary and sufficient
runtime implementation} of the guarantee.  Without it, the BX3 Intent Layer
is a specification without an execution path.

This relationship is the basis for regulatory defensibility.  An auditor who
accepts the BX3 Intent Layer's formal guarantee as a design commitment can
verify its enforcement by querying the Forensic Ledger for any trigger whose
\texttt{triggerCondition} is \texttt{accountability\_boundary}, confirming that
its \texttt{propagationChain} terminates at the declared \texttt{humanRoot}, and
verifying the corresponding \texttt{RESOLUTION\_RECORDED} entry.  The existence
of a queryable, cryptographically integrity-verified record is itself evidence
of the guarantee's operationalisation --- not merely its declaration.

The protocol thus closes the loop between BX3's abstract layer model and its
operational reality: the same formal structure that defines what each layer does
also defines how accountability is enforced when layers interact in ways that
violate the declared boundaries.  The Bailout Protocol is the point at which
BX3's architecture becomes accountable, not merely describable.

% ---------------------------------------------------------------
% End of Section 4
% ---------------------------------------------------------------